% Source: https://artofproblemsolving.com/community/c3h2479279_circleintersection_problem
% Stable identity: aops-topic-2479279-q01
\textbf{(a)} Substitute $x=2y-3$:
\[(2y-3)^2-3y^2+5(2y-3)=6-3y\iff4y^2-12y+9-3y^2+10y-15-6+3y=0\iff y^2+y-12=0.\]
So $(y+4)(y-3)=0$: $y=3$ gives $x=3$, and $y=-4$ gives $x=-11$. The points are $\boxed{A=(3,3),\ B=(-11,-4)}$.

\textbf{(b)} The midpoint of $AB$ is $M=\left(-4,-\dfrac12\right)$ and the slope of $AB$ is $\dfrac{3-(-4)}{3-(-11)}=\dfrac12$. The perpendicular bisector passes through $M$ with slope $-2$:
\[y+\dfrac12=-2(x+4)\iff\boxed{y=-2x-\dfrac{17}{2}}\quad(\text{i.e. }4x+2y+17=0).\]

\textbf{(c)} Completing the squares,
\[x^2+5x-3\left(y^2-y\right)=6\iff\left(x+\dfrac52\right)^2-3\left(y-\dfrac12\right)^2=6+\dfrac{25}{4}-\dfrac34=\dfrac{23}{2}.\]
The squares appear with opposite signs, so this is not a circle (nor a line) but a \textbf{hyperbola} with centre $\left(-\dfrac52,\dfrac12\right)$. The line meets both of its branches.

\begin{asy}
usepackage("amsmath");
size(7cm);
defaultpen(fontsize(10pt));
import graph;
real hp(real y){return -2.5+sqrt(11.5+3*(y-0.5)^2);}
real hm(real y){return -2.5-sqrt(11.5+3*(y-0.5)^2);}
draw((-14,0)--(9,0),Arrow(6)); draw((0,-6.5)--(0,7),Arrow(6));
label("$x$",(9,0),E); label("$y$",(0,7),N);
path b1=graph(hp,-5.6,6.6), b2=graph(hm,-4.7,5.6);
draw(reflect((0,0),(1,1))*b1,blue+linewidth(0.9));
draw(reflect((0,0),(1,1))*b2,blue+linewidth(0.9));
draw((-14,-5.5)--(9,6),gray+linewidth(0.8));
pair A=(3,3), B=(-11,-4), M=(-4,-0.5);
draw(A--B,red+linewidth(1));
draw((-7,5.5)--(-1.25,-6),dashed+heavygreen+linewidth(0.8));
pair u=unit(A-M), v=unit((-7,5.5)-M);
draw(M+0.8*u--M+0.8*u+0.8*v--M+0.8*v);
dot(A);dot(B);dot(M);
label("$A(3,3)$",A,SE); label("$B(-11,-4)$",B,SE); label("$M$",M,S);
label("$x=2y-3$",(-14,-5.5),S,gray);
label("$y=-2x-\dfrac{17}{2}$",(-7,5.5),N,heavygreen);
\end{asy}
